\documentclass[12pt, a4paper]{ctexart}
\usepackage[margin=2cm]{geometry}
\usepackage{amsmath, amssymb}
\usepackage{enumitem}
\usepackage{fancyhdr}
\usepackage{xcolor}
\usepackage{tikz}
\usepackage{graphicx}
\usepackage{tabularx}
\usepackage{multicol}

\setCJKmainfont{Hiragino Sans GB}
\setCJKsansfont{Heiti SC}

\pagestyle{fancy}
\fancyhf{}
\fancyhead[L]{培优专题·几何辅助线技巧}
\fancyhead[R]{测试卷}
\fancyfoot[C]{\thepage}
\renewcommand{\headrulewidth}{0.4pt}

\begin{document}

\begin{center}
{\Large\bfseries 培优专题·几何辅助线技巧 测试卷}\\[0.5em]
{\normalsize 满分150分\quad 考试时间120分钟}
\end{center}

\vspace{0.5em}

\noindent\textbf{一、单项选择题（本大题共8小题，每小题5分，共40分。每小题只有一个选项符合题意。）}

\begin{enumerate}[label=\Alph*.]

\item 在$\triangle ABC$中，$AB=5$，$AC=7$，$\angle BAC=60°$，$D$为$BC$的中点，则$AD$的长度为（\quad）
\begin{enumerate}[label=]
\item $\dfrac{\sqrt{19}}{2}$ \qquad \item $\dfrac{\sqrt{39}}{2}$ \qquad \item $\dfrac{\sqrt{37}}{2}$ \qquad \item $\dfrac{\sqrt{43}}{2}$
\end{enumerate}

\item 在$\triangle ABC$中，$D$是$BC$上一点，$BD:DC=3:1$，$E$是$AD$的中点，$BE$的延长线交$AC$于$F$，则$AF:FC$等于（\quad）
\begin{enumerate}[label=]
\item $1:2$ \qquad \item $2:3$ \qquad \item $3:4$ \qquad \item $1:3$
\end{enumerate}

\item 在$\triangle ABC$中，$\angle BAC=120°$，$AB=4$，$AC=3$，$D$为$BC$上一点，$AD$平分$\angle BAC$，则$AD$的长度为（\quad）
\begin{enumerate}[label=]
\item $\dfrac{12\sqrt{37}}{37}$ \qquad \item $\dfrac{12}{7}$ \qquad \item $\dfrac{12\sqrt{7}}{7}$ \qquad \item $\dfrac{6\sqrt{7}}{7}$
\end{enumerate}

\item 已知正方形$ABCD$的边长为$1$，$P$为正方形内一点，满足$PA^2+PC^2=PB^2+PD^2$，则$P$的轨迹为（\quad）
\begin{enumerate}[label=]
\item 一条线段 \qquad \item 一条直线 \qquad \item 一个圆 \qquad \item 两条线段
\end{enumerate}

\item 在$\triangle ABC$中，$BC=6$，$AC=8$，$\angle C=90°$，$I$为内心，$r$为内切圆半径，则$AI$的长度为（\quad）
\begin{enumerate}[label=]
\item $\sqrt{5}$ \qquad \item $2\sqrt{2}$ \qquad \item $\sqrt{6}$ \qquad \item $2\sqrt{5}$
\end{enumerate}

\item 在平行四边形$ABCD$中，$AB=8$，$BC=12$，$\angle ABC=60°$，$E$为$CD$的中点，则$AE$与$BD$的交点$P$满足$BP:PD$等于（\quad）
\begin{enumerate}[label=]
\item $2:1$ \qquad \item $3:1$ \qquad \item $3:2$ \qquad \item $5:2$
\end{enumerate}

\item 圆$O$的半径为$5$，弦$AB=8$，弦$CD=6$，$AB \parallel CD$，则$AB$与$CD$之间的距离为（\quad）
\begin{enumerate}[label=]
\item $1$ \qquad \item $7$ \qquad \item $1$或$7$ \qquad \item $2$或$6$
\end{enumerate}

\item 在$\triangle ABC$中，$G$为重心，$AD$为中线，$P$为$AD$上一点满足$AP:PD=3:1$，则$\triangle PBC$的面积与$\triangle ABC$的面积之比为（\quad）
\begin{enumerate}[label=]
\item $\dfrac{1}{2}$ \qquad \item $\dfrac{1}{3}$ \qquad \item $\dfrac{1}{4}$ \qquad \item $\dfrac{3}{8}$
\end{enumerate}

\end{enumerate}

\vspace{0.5em}

\noindent\textbf{二、多项选择题（本大题共4小题，每小题5分，共20分。全部选对得5分，部分选对得2分，有选错的得0分。）}

\begin{enumerate}[label=\arabic*.]

\item 在$\triangle ABC$中，$AB=AC$，$D$为$BC$的中点，$E$为$AD$上一点，下列结论正确的是（\quad）
\begin{enumerate}[label=]
\item $AD$是$BC$的垂直平分线
\item 若$BE \perp AC$，则$BD^2 = DE \cdot DA$
\item 若$\angle BEC = 90°$，则$BE$是$AC$边上的中线
\item $\triangle ABF$的面积$\leq \triangle ABD$的面积（$F$为$BC$上任意一点）
\end{enumerate}

\item 关于四边形的辅助线，下列说法正确的是（\quad）
\begin{enumerate}[label=]
\item 对角互补的四边形可以内接于圆
\item 连接梯形两条对角线中点的线段等于两底差的一半
\item 对角线互相垂直的四边形的面积等于对角线乘积的一半
\item 菱形的内切圆半径等于边长乘以一个角的正弦值
\end{enumerate}

\item 在$\triangle ABC$中，$H$为垂心，$O$为外心，$R$为外接圆半径，则下列关系正确的是（\quad）
\begin{enumerate}[label=]
\item $OH^2 = 9R^2 - (a^2+b^2+c^2)$
\item 若$\triangle ABC$为锐角三角形，则$H$在$\triangle ABC$内部
\item $AH = 2R\cos A$
\item $OH$的中点为九点圆的圆心
\end{enumerate}

\item 如图，$\triangle ABC$的内切圆与三边分别切于$D$、$E$、$F$，下列结论正确的是（\quad）
\begin{enumerate}[label=]
\item $AE = AF = \dfrac{b+c-a}{2}$
\item $\triangle DEF$的外接圆半径为$\dfrac{r}{2\sin\frac{A}{2}}$
\item $\triangle ABC$的面积$= r \cdot s$，其中$s$为半周长
\item $\angle EDF = 90° - \dfrac{A}{2}$
\end{enumerate}

\end{enumerate}

\vspace{0.5em}

\noindent\textbf{三、填空题（本大题共4小题，每小题5分，共20分。）}

\begin{enumerate}[label=\arabic*.]

\item 在$\triangle ABC$中，$AB=13$，$BC=14$，$AC=15$，$D$为$BC$的中点，则$AD=$ \underline{\hspace{3cm}}。

\item 在平行四边形$ABCD$中，$AB=6$，$AD=10$，$\angle BAD=120°$，$E$为$CD$的中点，连接$AE$、$BE$，则$\triangle ABE$的面积为 \underline{\hspace{3cm}}。

\item 已知$P$为$\triangle ABC$内一点，$\angle PAB = \angle PBC = \angle PCA = \alpha$，若$\triangle ABC$为边长为$1$的等边三角形，则$\alpha =$ \underline{\hspace{3cm}}°。

\item 在$\triangle ABC$中，$AB=10$，$AC=17$，$BC=21$，$AD$为$BC$边上的高，$M$为$AC$的中点，则$MD=$ \underline{\hspace{3cm}}。

\end{enumerate}

\vspace{0.5em}

\noindent\textbf{四、解答题（本大题共6小题，共70分。解答应写出文字说明、证明过程或演算步骤。）}

\begin{enumerate}[label=\arabic*.]

\item （本小题满分10分）\\
在$\triangle ABC$中，$\angle BAC = 90°$，$AB = AC$，$D$为$BC$上一点，$E$为$AB$上一点，满足$BD = AE$。过$D$作$DF \perp BC$交$AC$于$F$。\\
（1）求证：$\triangle AEF \cong \triangle BDE$；\\
（2）若$AB = 2$，$BD = \dfrac{1}{2}$，求$EF$的长度。
\vspace{6cm}

\item （本小题满分12分）\\
在$\triangle ABC$中，$AB = c$，$BC = a$，$CA = b$，$I$为内心，$r$为内切圆半径。\\
（1）求证：$S_{\triangle ABC} = \dfrac{1}{2}(a+b+c)r$；\\
（2）若$a = 13$，$b = 14$，$c = 15$，求内切圆半径$r$及$I$到顶点$A$的距离。
\vspace{6cm}

\item （本小题满分12分）\\
在四边形$ABCD$中，$AB = BC = CD = a$，$\angle ABC = 120°$，$\angle BCD = 60°$。\\
（1）证明：$A$、$B$、$C$、$D$四点共圆；\\
（2）求该四边形外接圆的半径。
\vspace{6cm}

\item （本小题满分12分）\\
在$\triangle ABC$中，$D$为$BC$上一点，$E$为$AD$上一点，满足$BD = 2DC$，$AE = 3ED$。$BE$的延长线交$AC$于$F$，$CE$的延长线交$AB$于$G$。\\
（1）求$\dfrac{AF}{FC}$的值；\\
（2）求$\dfrac{BG}{GA}$的值；\\
（3）求$\dfrac{S_{\triangle BEF}}{S_{\triangle ABC}}$的值。
\vspace{6cm}

\item （本小题满分12分）\\
已知$\odot O$的半径为$R$，$AB$为直径，$C$为圆上一点（$C$不与$A$、$B$重合），过$C$作$\odot O$的切线$l$，过$A$作$l$的垂线交$l$于$D$，过$B$作$l$的垂线交$l$于$E$。\\
（1）求证：$AD + BE = AB$；\\
（2）设$AC = m$，$BC = n$，求$DE$的长度（用$m$、$n$表示）。
\vspace{6cm}

\item （本小题满分12分）\\
在$\triangle ABC$中，$H$为垂心，$D$为$BC$的中点，$AH$的延长线交外接圆于$A'$。\\
（1）求证：$HD = DA'$；\\
（2）设$BC = a$，$AC = b$，$AB = c$，外接圆半径为$R$，求证：$AH = 2R\cos A$；\\
（3）若$BC = 10$，$AC = 7$，$AB = 5$，求$AH$的长度。
\vspace{6cm}

\end{enumerate}

\end{document}
